\documentclass[11pt]{article}
\usepackage[margin=1in]{geometry}
\usepackage{parskip}
\usepackage{hyperref}
\hypersetup{colorlinks=true, urlcolor=blue, linkcolor=black}
\sloppy

\begin{document}

\noindent Dear Editor,

\vspace{1em}

\noindent We would like to submit our manuscript, ``Adversarial Vulnerability of Machine
Learning GNSS Spoofing Detectors to Transmitted Spoofing Signals,'' for consideration as an
Original Article in \emph{Satellite Navigation}.

\vspace{1em}
\noindent\textbf{Why this manuscript belongs in \emph{Satellite Navigation}}

\vspace{0.5em}

\noindent Machine learning spoofing detectors are moving from research prototypes toward
integration in receivers, and the accuracies reported for them are high. Those accuracies
are obtained against fixed corpora of recorded attacks, which do not respond to the
detector being tested, while a deployed receiver faces an adversary that chooses what to
transmit. Our manuscript evaluates fourteen detectors, eight classical and six deep,
against attacks that are radio signals selected by an adversary. A counterfeit GPS L1 C/A
waveform is synthesised at sample level, summed with authentic recordings from the Texas
Spoofing Test Battery and tracked by an open software receiver, so the observables each
detector reads are produced by receiver processing of a transmittable signal. The generator
is validated by reproducing a recorded attack whose construction is published, and every
detector is held at one common operating point with the false alarm rate it pays reported
beside every result.

The central finding is that clean detection performance is a poor predictor of resistance
to transmitted attacks, and that near the top of the ranking the two orderings run close to
inverted: the detector with the highest clean $F_1$ score and the lowest false alarm rate
is evaded by 55 of 126 attacks, while the detectors evading none pay false alarm rates of
46 to 74 percent on authentic data. A missed attack is followed through the navigation
stage and shown to drive the receiver clock at 48.89 nanoseconds per second against a
commanded 48.87, so detection failure is not confined to the classifier. This speaks
directly to resilient positioning, navigation and timing, a core scope of the journal, and
to the growing literature on learning based GNSS spoofing detection that the journal has
published.

\vspace{1em}
\noindent\textbf{Policy issues}

\vspace{0.5em}

\noindent We are not aware of any policy issues relevant to this submission.

\vspace{1em}
\noindent\textbf{Competing interests}

\vspace{0.5em}

\noindent The authors declare that they have no competing interests.

\vspace{1em}
\noindent\textbf{Prior publication and authorship}

\vspace{0.5em}

\noindent This manuscript has not been published and is not under consideration for
publication elsewhere, in whole or in part. All authors have seen and approved the
manuscript in its current form and agree to its submission to \emph{Satellite Navigation}.

\vspace{1em}
\noindent\textbf{Special issue}

\vspace{0.5em}

\noindent This submission is not intended for a special issue.

\vspace{1em}

\noindent Thank you for considering our manuscript. We look forward to your response.

\vspace{1.5em}

\noindent Sincerely,

\vspace{1em}

\noindent Wasiu Akande Ahmed and Joel Segun Ojerinde (corresponding authors) \\
On behalf of all authors \\
Hangzhou International Innovation Institute, Beihang University,
Hangzhou 311115, China \\
\href{mailto:ahmedwasiu24@buaa.edu.cn}{ahmedwasiu24@buaa.edu.cn};
\href{mailto:joelojerinde@buaa.edu.cn}{joelojerinde@buaa.edu.cn}

\end{document}
